\section{成对比较能定出什么}
\label{sec:17-Constrained-Guarantees/05-Learning-from-Comparisons/01}

标注平台上跑着一个很朴素的任务：屏幕左右各放一段模型输出，标注员点一下更好的那个。三周下来收了四万条比较。现在要把这四万条``左边好''与``右边好''变成一个分数函数，好让它能当训练信号用。

有人问过为什么不省事一点，直接让标注员给每段输出打一个一到十分。试过，结果不能用。同一批文本发给两组人，两组的平均分差了一分半；把同一批文本隔一周再发给同一组人，同一个人的打分也漂了将近一分。而成对比较的重测一致率能到八成以上。这是一个可以重复观察到的经验规律，不是定理：\textbf{人对``哪个更好''的判断比对``值几分''的判断稳定得多}。绝对分数需要一把每个人心里各不相同的尺子，比较不需要。

于是问题变成：一堆两两比较的结果，究竟能把这把尺子定到什么程度？答案不是``数据够多就能定''，而是有几处结构性的空缺，数据量补不上。

\subsection{把比较写成一个概率模型}

给每个物品配一个实数分数 \(s_i\)，假设
\[
P(i\succ j)=\sigma(s_i-s_j),\qquad \sigma(u)=\frac{1}{1+e^{-u}} .
\]
这就是 Bradley--Terry 模型。它的形状与\hyperref[sec:13-Machine-Learning/02-Linear-Models/03]{``从二分类到 Softmax 分类''}里那个 logistic 模型完全一致，只是这里的``特征''是一个差分：把比较 \((i,j)\) 编码成一个只有第 \(i\) 位为 \(+1\)、第 \(j\) 位为 \(-1\) 的向量 \(e_i-e_j\)，则 \(P(i\succ j)=\sigma\bigl((e_i-e_j)^{\trans}s\bigr)\)。

给定 \(N\) 条独立比较，对数似然是
\[
\ell(s)=\sum_{(i,j)\in\mathcal{D}}\log\sigma(s_i-s_j),
\]
其中 \(\mathcal{D}\) 记录每条比较的胜方在前。这个函数对 \(s\) 是凹的——\(\log\sigma\) 是凹函数，凹函数与线性映射的复合仍然凹，求和保持凹性，这条路径在\hyperref[sec:09-Convex-Optimization/03-Convex-Functions/01]{``弦在图像上方意味着什么''}里已经走过。所以最大化 \(\ell\) 是一个凸问题，没有局部极大的困扰，而\hyperref[sec:11-Mathematical-Statistics/02-Point-Estimation/03]{``最大似然选择最能解释数据的参数''}给出的那套渐近性质原则上都适用。

原则上。下面三件事说明这个``原则上''前面要加几个限定。

\subsection{分数只由差值决定，所以整体平移不可见}

把所有 \(s_i\) 同时加上一个常数 \(c\)，每一个差 \(s_i-s_j\) 不变，于是每一条比较的概率不变，似然值一模一样。用矩阵语言说：设计矩阵的每一行都是 \(e_i-e_j\) 的形式，行和为零，所以全一向量 \(\mathbf{1}\) 落在设计矩阵的零空间里——这正是\hyperref[sec:03-Linear-Algebra/02-Vector-Spaces/03]{``矩阵的四个基本子空间：输入、输出与丢失的信息''}讲的那个方向：沿它移动，输出没有任何变化。

这不是``估计得不够准''。再来四百万条比较，似然沿 \(\mathbf{1}\) 方向仍然是平的，Fisher 信息矩阵在这个方向上的特征值恒为零。\hyperref[sec:11-Mathematical-Statistics/03-Sufficient-Statistics-and-Information/02]{``Score 与 Fisher 信息衡量局部可辨认性''}把这种情形叫做局部不可辨认，\hyperref[sec:11-Mathematical-Statistics/03-Sufficient-Statistics-and-Information/03]{``Cramér--Rao 界给无偏估计设下限''}的界在这个方向上退化为无穷——数据本身就问不出这个量。

要得到唯一解，必须外加一个约束，最常见的两种是把某一项固定为零（选一个参照物品），或者加一项 \(\lambda\norm{s}^2\) 的正则把解拉向原点，做法与\hyperref[sec:13-Machine-Learning/02-Linear-Models/02]{``正则化从 Ridge 到 Lasso''}里对付共线性时一样。两种做法都能给出唯一的数值解，但都改变不了一件事：\textbf{那个数值的绝对大小是约束造出来的，不是数据里的信息}。\(s_i=3.2\) 这句话本身不携带含义，只有 \(s_i-s_j\) 才携带。下一节会看到，这条限制一路传到奖励模型的用法上。

\subsection{比较图不连通，团与团之间就无法比较}

把物品当作顶点，每条比较当作一条边，得到一张比较图。这张图的连通性直接决定了能定出什么。

若图分成两个互不相连的部分 \(V_1\) 与 \(V_2\)，那么把 \(V_1\) 里所有分数加上 \(c_1\)、\(V_2\) 里所有分数加上 \(c_2\)，所有出现在数据里的差都不变，似然不变。于是自由度不止一个平移，而是每个连通分支各有一个。\hyperref[sec:04-Discrete-Mathematics/03-Graph-Theory/02]{``路径、环与连通性''}里那个判据在这里直接就是可辨识性判据：\textbf{连通分支的个数等于不可辨识的平移自由度个数}。

这一条最容易被误读成数据量问题。它不是。在 \(V_1\) 内部再加一百万条比较，\(V_1\) 内部的分数会估得更准，\(V_1\) 与 \(V_2\) 之间的关系仍然一无所知——因为从来没有一条边跨过去。补救办法只有一个：真的去采集跨团的比较。这在实践中很容易被忽略，因为采样规则常常是``在同一个 prompt 下比较两个回答''，而不同 prompt 之间从不比较，于是比较图天然按 prompt 分块。若还要把不同 prompt 上的分数放在一起用，那就是在使用一个从未被数据约束过的量。

严格地说，连通只是必要条件。若某个分支内部存在一组物品从未输过，似然的最大值会跑到无穷远，参数估计不存在——这与 logistic 回归在可完全分离数据上的表现是同一件事，正则化项同时也在防这个。

\begin{center}
\begin{tikzpicture}[x=1cm,y=1cm,font=\small,>=stealth,
  itm/.style={draw,circle,minimum size=0.52cm,inner sep=0pt,fill=blue!8}]
  % 左：连通
  \begin{scope}
    \node[itm] (n1) at (0.2,1.5) {};
    \node[itm] (n2) at (1.9,1.8) {};
    \node[itm] (n3) at (3.2,0.8) {};
    \node[itm] (n4) at (0.5,0.2) {};
    \node[itm] (n5) at (2.2,-0.1) {};
    \draw[black!60] (n1) -- (n2);
    \draw[black!60] (n2) -- (n3);
    \draw[black!60] (n1) -- (n4);
    \draw[black!60] (n4) -- (n5);
    \draw[black!60] (n2) -- (n5);
    \draw[black!60] (n3) -- (n5);
    \node[align=center,font=\footnotesize] at (1.7,-1.0)
      {连通：分数可定\\只差一个整体平移};
  \end{scope}
  % 右：不连通
  \begin{scope}[xshift=5.9cm]
    \node[itm] (a1) at (0.2,1.6) {};
    \node[itm] (a2) at (1.4,1.8) {};
    \node[itm] (a3) at (0.8,0.5) {};
    \draw[black!60] (a1) -- (a2);
    \draw[black!60] (a2) -- (a3);
    \draw[black!60] (a1) -- (a3);
    \node[itm] (b1) at (2.7,1.6) {};
    \node[itm] (b2) at (3.9,1.4) {};
    \node[itm] (b3) at (3.2,0.4) {};
    \draw[black!60] (b1) -- (b2);
    \draw[black!60] (b2) -- (b3);
    \draw[black!60] (b1) -- (b3);
    \draw[black!45,dashed] (2.05,-0.15) -- (2.05,2.2);
    \node[align=center,font=\footnotesize] at (2.0,-1.0)
      {不连通：团内可定\\团间关系无法确定};
  \end{scope}
  % 下：非传递循环
  \begin{scope}[yshift=-3.5cm,xshift=2.2cm]
    \node[itm] (t1) at (0,0) {\footnotesize\(A\)};
    \node[itm] (t2) at (2.2,0) {\footnotesize\(B\)};
    \node[itm] (t3) at (1.1,1.5) {\footnotesize\(C\)};
    \draw[->,black!70,thick] (t1) -- (t2);
    \draw[->,black!70,thick] (t2) -- (t3);
    \draw[->,black!70,thick] (t3) -- (t1);
    \node[anchor=west,align=left,font=\footnotesize] at (3.0,0.75)
      {非传递的循环：\(A\) 多数胜 \(B\)，\(B\) 多数胜 \(C\)，\(C\) 多数胜 \(A\)。\\
       任何一维分数都排不出这种关系，残差里会留下系统结构。};
  \end{scope}
\end{tikzpicture}

\emph{比较数据能定出什么，由这张图的形状决定：连通性决定有几个平移自由度，环的方向决定一维分数够不够用。数据量在这两件事上都不起作用。}
\end{center}

\subsection{一维分数假定了传递性}

Bradley--Terry 里，\(s_i>s_j\) 当且仅当 \(P(i\succ j)>1/2\)。既然 \(s\) 是一组实数，实数的序是全序，模型就隐含地断言了整体偏好可以排成一条线：若 \(A\) 多数胜 \(B\)、\(B\) 多数胜 \(C\)，则必有 \(A\) 多数胜 \(C\)。

真实的成对偏好不保证有这个性质。标注指南里同时写着``要准确''与``要简洁''，而三段输出分别在不同维度上占优时，两两比较的多数结果完全可以形成 \(A\succ B\succ C\succ A\) 的循环。这不需要标注员不理性，只需要他们在不同的比较对上把注意力放在不同的维度——每一次判断都自洽，合起来不传递。上图底部那个三角就是这种情形的最小形态。

Bradley--Terry 遇到循环时不会报错，它会给出一组分数，让三个概率都尽量靠近 \(1/2\)，从而把循环``抹平''成三个近似平局。这时模型的预测精度掉下来，而掉下来的方式有特征：残差不是随机分布的，它在循环涉及的那几条边上系统性地同号。检查残差有没有这种结构，是判断一维分数是否够用的一个实际手段。若确实存在稳定的循环，补救方向不是加数据，而是承认偏好至少是多维的——把评价拆成几个维度各自建模，或者干脆放弃全序。

\subsection{这一节定住了什么}

比较数据能定出的是差值，不是水平；能定出的范围以比较图的连通分支为界；而把差值压进一个实数分数这件事本身，已经假设了偏好可传递。三条限制里，前两条是模型的代数结构和图结构给的，可以事先检查；第三条是关于世界的假设，只能事后从残差里查。

下一节把 \(s_i\) 换成一个以输入为条件的函数 \(r_\theta(x,y)\)，把上面这个对数似然直接当成训练损失。换掉的只是参数化方式，三条限制一条都不会消失——它们会以更不容易被注意到的形式重新出现。
